%%%%%%%%%%%%%%%%%%%%%%%%%%%%%%%%%%%%%%%%%%%%%%%%%%%%%%%%%%%%%%%%%%%%%%
% Lớp 9 | Chương IV. Hệ thức lượng trong tam giác vuông | Tỉ số lượng giác của góc nhọn
% Dạng 9H4?11-4: Rút gọn, chứng minh biểu thức lượng giác
%%%%%%%%%%%%%%%%%%%%%%%%%%%%%%%%%%%%%%%%%%%%%%%%%%%%%%%%%%%%%%%%%%%%%%


% nguồn: dan_tu_AI (nhập 17/11/2026)
\begin{ex}%[9H4H11-4]%[ma:66dc3a44]
	Giá trị của biểu thức $A=\sin^2 25^\circ+\sin^2 65^\circ$ là
	\choice
	{\True $1$}
	{$2$}
	{$0$}
	{$\dfrac12$}
	\loigiai{
		Vì $25^\circ+65^\circ=90^\circ$ nên $\sin 65^\circ=\cos 25^\circ$.\\
		Do đó $A=\sin^2 25^\circ+\cos^2 25^\circ=1$.
	}
\end{ex}

% nguồn: dan_tu_AI (nhập 17/11/2026)
% nguồn: dan_tu_AI (nhập 17/11/2026)
% nguồn: dan_tu_AI (nhập 17/11/2026)
